


\section{Disclosure selection and the observation likelihood}
Imagine equal prior probabilities of a good and bad type. If good types
disclose with probability $0.9$ and bad types with probability $0.1$, then
disclosure raises the probability of a good type to $0.9$, while silence
reduces it to $0.1$. Treating a missing disclosure as no information would
retain $0.5$ in both cases. The discrepancy comes from omitting the selection
mechanism, not from numerical approximation. These probabilities are an
authored example, not estimates of actual issuer disclosure behaviour.

Let $Y$ denote potentially observable information, $M$ the pattern of
observed components, $f_\theta(Y)$ its model and $g_\phi(M\mid Y)$ the
publication or missingness mechanism. The likelihood of the realised data is
\begin{equation}
L(\theta,\phi)=\int
f_\theta(y_{\rm obs},y_{\rm mis})
g_\phi(m\mid y_{\rm obs},y_{\rm mis})\,dy_{\rm mis}.
\end{equation}
Here the realised pattern $m$ fixes which coordinates are observed; the
integration is over the missing coordinates only. Discrete missing components
use a sum instead. If, for the realised observation, $g_\phi(m\mid y)$ is constant over
all compatible missing components, it factors out of this integral. Together
with distinct parameter spaces this yields the relevant likelihood
proportionality; a Bayesian factorisation also needs the appropriate prior
independence. Seaman and coauthors distinguish these realised-MAR conditions
from the everywhere-MAR conditions relevant to repeated-sampling comparisons
\cite[Theorems~1--3, Section~5.4 and Appendix]{litcloseseaman}. These are sufficient
conditions for specified inferential equivalences. They do not demonstrate
that corporate missingness satisfies them.

\section{When silence has an equilibrium interpretation}
A separate local illustration makes non-disclosure endogenous. Let value $V$
be uniform on $[0,1]$. A manager knows $V$ with probability $1/2$, independently
of value; otherwise the manager cannot disclose it. Informed disclosure is
truthful and costless, the manager maximizes the contemporaneous price, and
risk-neutral investors price at conditional expected value. Suppose an informed
manager discloses above cutoff $a$. Silence occurs with probability
$(1+a)/2$. Its unconditional value contribution is
$1/4+a^2/4$, giving silence price
\begin{equation}
P_{\rm silence}(a)=\frac{1+a^2}{2(1+a)}.
\end{equation}
An informed manager at the cutoff is indifferent, so $a=P_{\rm silence}(a)$.
Multiplying through and collecting terms gives $a^2+2a-1=0$, hence
$a=\sqrt2-1\simeq0.4142$. Types above that value disclose and those below
it prefer the silence price. The uninformed manager explains why silence
does not identify the worst type. This derivation is a local example; the
blocked Dye original is not claimed to have been reproduced.

Nor is favorable selection universal. Suijs models a firm that knows its
return but is uncertain how investors compare it with an outside opportunity
\cite[Sections~3--4, Proposition~2 and Appendices~A--B]{litclosesuijs}. Under the
stated conditions a profitable firm can disclose lower returns and withhold
higher returns. Multiple partial-disclosure equilibria are possible. The
implication for AccDSL is to represent alternative disclosure mechanisms and
test their consequences, rather than hard-code one direction of selection.

\section{What changes under an intervention}
A mechanical rate stress replaces an input while holding specified other
relations fixed. A conditional forecast computes outcomes among states where
an observed rate has a chosen value. A causal intervention changes the
mechanism generating the rate and then propagates the remaining structural
relations. These calculations can disagree.

For a finite counterexample, let $U$ be a fair binary variable and observe
$X=U$, $Y=U$. An alternative structural model has $X=U$, $Y=X$. Both give
exactly the same observed joint distribution. Under an intervention setting
$X=1$, the first model retains $\E[Y]=1/2$ while the second gives
$\E[Y]=1$. Observational agreement does not identify the intervention
response. A forecast labelled as causal therefore needs the structural
assumptions that distinguish these possibilities, or must retain the set of
responses. Source corrections and simulated observations also have different
time meanings: a hypothetical future disclosure updates an actor inside that
run, not the analyst's actual present evidence.
